\documentclass[10pt]{article}

\pdfinfoomitdate=1
\pdftrailerid{}
\pdfsuppressptexinfo=-1

\usepackage[T1]{fontenc}
\usepackage{lmodern}
\usepackage[margin=1in]{geometry}
\usepackage{microtype}
\usepackage{amsmath}
\usepackage{booktabs}
\usepackage{url}
\usepackage[hidelinks]{hyperref}

\hypersetup{pdftitle={The Quantized Mode of KodamaEmbed: Bit-Identical Centered INT8 BERT Embeddings on CPUs and NVIDIA GPUs},
            pdfauthor={Yuriy Peysakhov}}

\title{The Quantized Mode of KodamaEmbed:\\Bit-Identical Centered INT8 BERT Embeddings\\on CPUs and NVIDIA GPUs}
\author{Yuriy Peysakhov}
\date{}

\begin{document}
\maketitle

\begin{abstract}
KodamaEmbed is an open-source Rust engine for BERT and ModernBERT embedding
and reranking models with CPU, CUDA and Vulkan backends. Its quantized mode,
centered INT8, specifies one
numerical contract for the whole encoder: static per-token-class centers,
rotated activation quantization in 128-value blocks, GPTQ weight digits, exact
integer dot products accumulated through a fixed FP32 chain, and integer
attention with tabulated 14-bit softmax weights. Every reduction in the
contract is exact or has a prescribed order, so a sequence's vector depends
only on its tokens. The CPU and CUDA backends implement the contract with
identical vector bits for every input length and batch composition, which the
test suite checks at every supported length and with exhaustive checks of the
elementary functions. Version 0.1.0 runs the mode for E5-large-v1; its other
models run in float precision. On an
NVIDIA RTX 4000 SFF Ada Generation, KodamaEmbed in this mode is the fastest
engine at seven fixed query lengths from 16 to 512 tokens: the fastest of
ONNX Runtime,
TensorRT and PyTorch at FP16 takes 1.07 to 2.15 times as long. On six
efficiency cores of an Intel Core i5-13500, the faster of ONNX Runtime and
PyTorch at FP32 takes 1.80 to 2.38 times as long at those lengths.
Binary-relevance nDCG@10 differs
from FP32 by less than 0.002 on each of four dataset aggregates. All 13
individual retrieval tasks pass the quality gates.
\end{abstract}

\section{Introduction}

Retrieval systems embed a corpus once and embed queries continuously. Corpus
vectors must remain comparable with later query vectors, which may be computed
on another device or in a different batch. Floating-point kernels can sum the
same inputs in different orders, changing the last bits of a vector when a
runtime selects a different execution plan. For short queries, kernel launches
and synchronization also contribute to end-to-end latency.

KodamaEmbed addresses both with an exactly specified approximate encoder. The
centered INT8 mode replaces floating-point accumulation with exact integer
arithmetic wherever the computation allows, and fixes the order of every
remaining floating-point operation. Its CPU and CUDA implementations then
agree to the bit, and a vector does not depend on the batch in which it was
computed.
Corpus and query vectors may come from either device. The CUDA backend is
built for short inputs: a batch of up to 16 tokens can run every layer in one
cooperative launch.

This report makes four contributions:
\begin{itemize}
\item a numerical contract for BERT embedding encoders in which every
  reduction is exact or has a prescribed order (Section~\ref{sec:contract});
\item CPU and CUDA implementations of that contract with identical vector
  bits, and the tests that verify them (Section~\ref{sec:identity});
\item a CUDA execution design for short inputs (Section~\ref{sec:cuda});
\item a matched latency comparison against three FP16 runtimes on a GPU and
  two FP32 runtimes on a processor, and a retrieval-quality evaluation
  against FP32, with published receipts (Section~\ref{sec:evaluation}).
\end{itemize}

The centered mode approximates FP32 inference and accepts one exact
checkpoint, E5-large-v1~\cite{wang2022e5}; the performance results come from
one GPU and one processor. Section~\ref{sec:limitations} states the scope of
the results.

\section{Background and related work}

\paragraph{Embedding encoders.}
BERT~\cite{devlin2019} encoders produce sentence embeddings by pooling their
final hidden states; Sentence-BERT~\cite{reimers2019} established mean-pooled
siamese training, MiniLM~\cite{wang2020minilm} distills compact encoders, and
E5~\cite{wang2022e5} trains embeddings contrastively with query and passage
prefixes. KodamaEmbed serves these and other BERT encoders and
ModernBERT~\cite{warner2025} encoders, pooled and L2-normalized as each model
card prescribes.

\paragraph{Quantization.}
Post-training INT8 inference for transformers must handle activation
outliers~\cite{dettmers2022,xiao2023}. GPTQ~\cite{frantar2023} chooses weight
digits that compensate rounding error using a calibration Hessian.
Rotations spread outliers across channels before quantization;
QuaRot~\cite{ashkboos2024} applies Hadamard transforms for this purpose.
Centering is also established: HadaNorm~\cite{federici2025} removes per-channel
means before a Hadamard transform, and SageAttention~\cite{zhang2025sage}
and SageAttention2~\cite{zhang2025sage2} subtract mean components of attention
operands before low-bit quantization. REx~\cite{yvinec2023} expands a
quantized network in residual quantization errors. KodamaEmbed combines static
centers per token class, block Hadamard rotation, GPTQ digits and one residual
term, a second INT8 digit of the operand at two projections; its contribution
is the exactly specified arithmetic around these established techniques.

\paragraph{Reproducibility.}
Floating-point addition is not associative, so changing the order of a
parallel sum can change its result~\cite{demmel2015}.
Batch-invariant kernels remove the dependence of LLM outputs on batch
composition within one implementation~\cite{he2025}. KodamaEmbed requires more:
equal bits across a CPU and a GPU implementation. It obtains them by
accumulating in exact integers and fixing the order of every remaining
floating-point operation.

\section{The centered INT8 contract}
\label{sec:contract}

\subsection{Projections}

Each token belongs to one of three static classes: the first token, the last
token and the interior tokens. Calibration records each class's mean input to
every projection; its FP16 rounding is the class center $c_t$. A projection
computes
\begin{equation}
y = W(x - c_t) + W c_t + b,
\end{equation}
where the shared row $W c_t$ is computed once from the exact FP32 weights in
binary64 and rounded once. The operand $x$ is rounded to FP16, centered,
multiplied by fixed signs, rotated by a normalized Hadamard transform within
128-value blocks, and quantized to INT8 with one step per row and half block,
the 64 inputs of one part's dot: the half's largest magnitude divided by 127.
The weights receive the same signs
and rotation, one scale per output channel, and INT8 digits chosen by GPTQ on
the Hessian of the centered, rotated calibration operand. The GPTQ solve uses
Babai's nearest-plane form on a reversed Cholesky factor, and every sum in it
runs in a fixed order, so the digits do not depend on the thread count.

The product accumulates exact INT32 dots over 64-input parts. Each part's dot
converts to FP32 exactly and joins the row's sum by one fused multiply-add on
its part's step, in depth order. The chain restarts at each multiple of the
hidden width, and the feed-forward down projection adds its segments in a fixed
order.
The output-channel scale, the shared row, the bias and the residual or GELU
follow in a fixed order.

\paragraph{A second operand digit.}
The last layer's QKV projection and the up projection of the layer before it
keep a second INT8 digit of each rotated operand value. With $d_1$ the digit
on the half block's step $s$, the second digit rounds what the first leaves,
$\mathrm{fma}(-d_1, s, x)$, on the step $s/254$; the first leaves at most
$s/2$, so the second stays within 127. The product's chain runs on through the
second digits' parts over the same weight digits, each on its step $s/254$.
The weights, the centers and the calibrated tables do not change.

\paragraph{Deferred LayerNorm.}
For the LayerNorms before the feed-forward up projection and the later layers'
QKV projections, the encoder uses
\begin{equation}
W\bigl(\gamma \odot (z - \mu)\,r + \beta\bigr) + b
= r\bigl(W(\gamma \odot z) - \mu\,W\gamma\bigr) + W\beta + b,
\end{equation}
so the row's statistics $\mu$ and $r$ enter only the projection epilogue and
the raw row is quantized first.

\subsection{Integer attention}

Queries are quantized to INT8 per row. Keys are first centered on a static
key center, which shifts all of a query's scores by the same amount and so
leaves softmax unchanged, then quantized per row. A score is the exact integer
dot times the key's scale. The exponent of a score's distance from its row's
largest score, in units of $2^{-8}$ in $\log_2$, indexes a table of 14-bit
weights $\mathrm{RNE}(16383 \cdot 2^{-u/256})$. At a whole octave the entry
rounds the exact quotient $16383 / 2^k$, and every other entry lies far from
a rounding tie, so any faithful $\exp_2$ produces the same table. The
weights' sum is an exact integer.

Values are quantized to INT8 per sequence and channel, with a step taken over
the non-sink tokens. Each 14-bit weight splits into two 7-bit digits, so the
weighted sums of values use INT8 products and stay exact in INT32 up to 512
keys ($512 \cdot 16383 \cdot 127 < 2^{31}$). The first and last tokens draw the
most attention, so their values keep FP16 precision: their entries in the
quantized values are zero, and the last step adds their terms explicitly:
\begin{equation}
o = \bigl((p\,s + w_{\mathrm{first}} v_{\mathrm{first}})
      + w_{\mathrm{last}} v_{\mathrm{last}}\bigr)\cdot\frac{1}{T},
\end{equation}
where $p$ is the exact integer product sum, $s$ the value step and $T$ the
integer weight total.

\subsection{Element-wise operations and pooling}

GELU uses the Abramowitz and Stegun erf approximation~\cite{abramowitz1964}
(7.1.28) with a prescribed sequence of fused multiply-adds. LayerNorm
statistics sum in a fixed lane order. Mean pooling sums FP16 values as exact
integers in units of $2^{-24}$; a sequence of at most 512 values stays below
$2^{49}$ units, where every binary64 partial sum is exact. The mean is formed
by integer long division with a sticky bit, which rounds as the binary64
division would. The norm sums squares in a fixed order, and each output
equals the CPU's binary64 quotient rounded to FP32.

\subsection{Consequences}

Each sequence's vector is a deterministic function of its tokens and the
calibrated tables. It does not depend on the other sequences in a batch, the
thread count or the device. Calibration runs one FP32 forward pass over 128
passages shipped with the engine, over queries formed from each passage's
first 16 words, and over 1,024 queries of one to three vocabulary pieces drawn
by a fixed generator, then solves GPTQ per projection. Its tables are cached
under a
fingerprint that covers the model files, prefixes, truncation, numerical
sources and platform.

\section{Bit-identical implementations}
\label{sec:identity}

The CPU implementation uses AVX2, FMA and F16C, with AVX-VNNI for the integer
products. The CUDA implementation performs the same IEEE operations in the
same order. Explicit rounding intrinsics keep the compiler from contracting or
reordering floating-point operations, and fused multiply-adds appear exactly
where the CPU uses them. Work is split across threads only where order cannot
matter: exact integer sums, exact binary64 sums of FP16 values, and maxima.

On the GPU, a quotient is formed from a correctly rounded reciprocal, computed
without the approximate hardware reciprocal, and one fused multiply-add
correction. For the divisors that occur, integers below
$2^{21}$, and dividends of magnitude $2^{-60}$ to $2^{60}$, every intermediate
stays normal; an exhaustive check of every dividend in $[1, 2)$ against every
divisor in $[2^{20}, 2^{21})$ therefore covers the range by scaling.

The test suite compares the CUDA and CPU encoders bit for bit at every input
length from 1 to 512 tokens and on random batches, including the thresholds
between execution plans and storage reuse. Exhaustive
checks cover GELU on all $2^{32}$ floats, the quotient range above, the
attention-weight index and quantization steps on every float, and the integer
conversion of every FP16 value. The reciprocal is checked against the
division on every normal float below $2^{126}$.

Calibration is deterministic as well: the GPTQ solve and every calibration
sum run in a fixed order. In the release measurement, CUDA vectors for 26
real-data files, 27,541 rows of SciFact, NFCorpus, FiQA and LoCoMo passages
and queries, equal the CPU's vectors byte for byte.

\section{CUDA execution}
\label{sec:cuda}

The encoder packs variable-length sequences without padding. Batches of more
than 16 tokens run per-layer kernels. Projections use INT8 tensor-core
products with tile shapes chosen by row count, and their epilogues fuse the
residual, GELU and the next projection's quantization, or the packing of
attention operands. Batches of up to 1,024 tokens replay cached CUDA graphs.

Batches of up to 16 tokens can run a persistent plan that executes every
layer in one cooperative launch. Each multiprocessor holds one resident block.
Weights stream through it with evict-first loads, so they do not displace the
activations in L2, while worker warps prepare row-local values; counted
hand-offs order shared inputs between blocks. Shared-memory
stream slots have an explicit completion contract: each participating lane
acquires the barrier phase before reading the slot, and warp participation
reconverges before execution proceeds. A plan runs only after passing
geometry, multiprocessor-layout, cooperative-launch and residency checks;
otherwise the batch uses the per-layer route with the same arithmetic.

\section{Evaluation}
\label{sec:evaluation}

All measurements use KodamaEmbed 0.1.0 on one system: Linux 6.8, an Intel
Core i5-13500 and an NVIDIA RTX 4000 SFF Ada Generation (driver 595.112.07)
under its 70\,W power limit. The model is E5-large-v1 (24 layers, hidden
width 1024) at a pinned revision, with mean pooling and L2 normalization.

\subsection{Latency}

Each call takes one query text and returns its normalized FP32 embedding on
the host, so tokenization, transfer, inference, pooling and normalization are
timed together. Every engine receives the same text, prefix, truncation,
token IDs and segment IDs. Inputs are SciFact word prefixes whose encoded
lengths, including special tokens and E5's \texttt{query:~} prefix, are
exactly the lengths in Table~\ref{tab:latency}; at the model's maximum length
the text is truncated. KodamaEmbed runs centered INT8 through its public API.
On the GPU, ONNX Runtime~1.30.0 runs an FP16 graph with TF32 disabled and CUDA
graph replay; TensorRT~10.16.1.11 runs through ONNX Runtime's TensorRT
provider at FP16 with builder optimization level~5, a fixed shape per case and
CUDA graph replay; PyTorch~2.14.1 runs FP16 BERT with scaled-dot-product
attention and CUDA graphs. On the processor, ONNX Runtime and PyTorch run FP32
graphs, and every engine has six threads. The other engines pool and
normalize with FP32 operations.

One benchmark process runs at a time, in three rounds that reverse the
previous round's order. Each process settles for 30 seconds after loading,
warms up for at least 2 seconds and 20 calls, and records 51 samples of three
calls, keeping every returned output alive until its sample's clock stops.
Processes run at nice~19 on six efficiency cores. Before each GPU process,
three consecutive GPU samples five seconds apart must show at most
86\,\textdegree C, an enforced power limit of at least 70\,W, and no software
power cap or thermal slowdown; 100\,ms telemetry must cover the timed window,
and a low enforced limit or thermal slowdown during timing rejects the
process. The GPU also held another workload's idle compute context, and every
gate sample showed 0\% utilization. Each result is the median of three
process medians. Table~\ref{tab:latency} lists all 14 workloads.

\begin{table}[t]
\centering
\setlength{\tabcolsep}{5pt}
\caption{Single-query latency of E5-large-v1 in milliseconds, on the GPU
(top) and on six efficiency cores (bottom), each the median of three process
medians. The last column divides the fastest other engine's latency by
KodamaEmbed's, from unrounded values.}
\label{tab:latency}
\begin{tabular}{rrrrrr}
\toprule
Tokens & KodamaEmbed & ONNX Runtime & TensorRT & PyTorch & Fastest other \\
 & INT8 & FP16 & FP16 & FP16 & / KodamaEmbed \\
\midrule
16 & 1.325 & 2.910 & 2.846 & 3.189 & 2.15 \\
33 & 1.594 & 3.035 & 2.932 & 3.395 & 1.84 \\
64 & 1.814 & 3.170 & 3.088 & 3.641 & 1.70 \\
65 & 2.044 & 3.323 & 3.165 & 3.747 & 1.55 \\
128 & 2.595 & 3.851 & 3.476 & 4.022 & 1.34 \\
256 & 4.643 & 5.173 & 4.968 & 5.599 & 1.07 \\
512 & 7.777 & 9.635 & 8.415 & 9.897 & 1.08 \\
\bottomrule
\end{tabular}

\vspace{2ex}
\begin{tabular}{rrrrr}
\toprule
Tokens & KodamaEmbed & ONNX Runtime & PyTorch & Fastest other \\
 & INT8 & FP32 & FP32 & / KodamaEmbed \\
\midrule
16 & 23.7 & 56.4 & 87.2 & 2.38 \\
33 & 45.9 & 97.4 & 125.9 & 2.12 \\
64 & 85.4 & 173.1 & 293.1 & 2.03 \\
65 & 87.1 & 176.5 & 294.2 & 2.03 \\
128 & 183.7 & 337.6 & 446.4 & 1.84 \\
256 & 376.9 & 678.5 & 783.7 & 1.80 \\
512 & 771.8 & 1466.2 & 1474.6 & 1.90 \\
\bottomrule
\end{tabular}
\end{table}

At each of the seven fixed lengths, KodamaEmbed with centered INT8 is the
fastest engine on both devices: its slowest process median is below the
fastest process median
of every FP16 competitor on the GPU and every FP32 competitor on the
processor. The second-fastest engine on the GPU, TensorRT at each length,
takes 1.07 to 2.15 times as long as centered INT8; on the processor, ONNX
Runtime takes 1.80 to 2.38 times as long.
A second measurement runs one text of each length from 16 to 512 tokens in a
fixed order. A fresh session makes three passes without a separate warm-up.
On the later two passes, centered INT8 again has the lowest mean latency;
the second lowest, ONNX Runtime's, is 1.27 times it.

\subsection{Retrieval quality}

Centered INT8 approximates FP32, so its vectors are compared with FP32 CPU
vectors of the same data under fixed gates: mean cosine distance at most
0.001 for passages and queries, maximum L2 deviation at most 0.1, unit-norm
error at most $10^{-5}$, mean top-10 overlap at least 0.9, and losses of at
most 0.005 in nDCG@10 and 0.01 in recall@10 and hit@10. Table~\ref{tab:quality}
reports the evaluation on SciFact, NFCorpus and the first 10,000 passages of
FiQA from BEIR~\cite{thakur2021}, and on evidence-turn retrieval over
LoCoMo~\cite{maharana2024}. nDCG@10 uses binary relevance.

\begin{table}[t]
\centering
\caption{Centered INT8 against FP32 for E5-large-v1: exhaustive dot-product
retrieval.}
\label{tab:quality}
\begin{tabular}{lrrrrr}
\toprule
Dataset & Queries & \multicolumn{2}{c}{nDCG@10} & Query cos.\ dist. & Top-10 overlap \\
 & & FP32 & Centered & & \\
\midrule
SciFact & 300 & 0.7256 & 0.7251 & $2.4\times10^{-5}$ & 0.989 \\
NFCorpus & 323 & 0.3616 & 0.3616 & $2.1\times10^{-5}$ & 0.985 \\
FiQA & 243 & 0.5096 & 0.5076 & $2.1\times10^{-5}$ & 0.988 \\
LoCoMo & 1977 & 0.5506 & 0.5503 & $2.3\times10^{-5}$ & 0.989 \\
\bottomrule
\end{tabular}
\end{table}

All four dataset aggregates pass the quality gates, with binary-relevance
nDCG@10 differing from FP32 by less than 0.002 in each. All 13 individual
tasks also pass their gates: one each for SciFact, NFCorpus and FiQA, and ten
LoCoMo conversations. The largest individual nDCG@10 loss is 0.00401 on a
LoCoMo conversation, within the 0.005 loss limit. FiQA was held out during
design and calibration.

\section{Limitations}
\label{sec:limitations}

The centered mode approximates FP32 inference and accepts the one evaluated
checkpoint, verified by its file hashes. Calibration and CPU-side
preparation require an x86-64 host with AVX-VNNI, AVX2, FMA and F16C, and CUDA
execution requires compute capability 8.0 or later. All timings come from one
GPU under its 70\,W power limit and six efficiency cores of one processor, on
a server that also ran other workloads, and measure sequential single-query
latency; throughput, concurrent serving and cold starts are outside this
evaluation. In the fixed-length comparison, all three GPU competitors used
FP16 and CUDA graph replay. In the changing-length comparison, ONNX Runtime
and TensorRT used dynamic shapes without CUDA graph replay, and PyTorch ran
eagerly. Three of the four quality datasets were
used during development. Bit-exact comparisons on the tested device and
workloads do not establish freedom from data races on other devices or
workloads.

\section{Conclusion}

An approximate encoder can be exactly specified. By accumulating in
exact integers and fixing the order of every remaining floating-point
operation, KodamaEmbed makes centered INT8 BERT embeddings identical across its
CPU and CUDA backends and independent of batch composition, while serving
short queries at low latency. The source, its tests, the benchmark harness and
the receipts behind every measurement in this report are released with
KodamaEmbed under the Apache-2.0 license.

\begin{thebibliography}{99}
\small

\bibitem{abramowitz1964}
M.~Abramowitz and I.~A. Stegun, editors.
\newblock \emph{Handbook of Mathematical Functions with Formulas, Graphs, and
  Mathematical Tables}.
\newblock National Bureau of Standards, 1964.

\bibitem{ashkboos2024}
S.~Ashkboos, A.~Mohtashami, M.~L. Croci, B.~Li, P.~Cameron, M.~Jaggi,
  D.~Alistarh, T.~Hoefler, and J.~Hensman.
\newblock {QuaRot}: Outlier-free 4-bit inference in rotated {LLMs}.
\newblock In \emph{Advances in Neural Information Processing Systems}, 2024.
\newblock arXiv:2404.00456.

\bibitem{demmel2015}
J.~Demmel and H.~D. Nguyen.
\newblock Parallel reproducible summation.
\newblock \emph{IEEE Transactions on Computers}, 64(7):2060--2070, 2015.

\bibitem{dettmers2022}
T.~Dettmers, M.~Lewis, Y.~Belkada, and L.~Zettlemoyer.
\newblock {LLM.int8()}: 8-bit matrix multiplication for transformers at scale.
\newblock In \emph{Advances in Neural Information Processing Systems}, 2022.
\newblock arXiv:2208.07339.

\bibitem{devlin2019}
J.~Devlin, M.-W. Chang, K.~Lee, and K.~Toutanova.
\newblock {BERT}: Pre-training of deep bidirectional transformers for language
  understanding.
\newblock In \emph{Proceedings of NAACL-HLT}, 2019.
\newblock arXiv:1810.04805.

\bibitem{federici2025}
M.~Federici, R.~Del~Chiaro, B.~van Breugel, P.~Whatmough, and M.~Nagel.
\newblock {HadaNorm}: Diffusion transformer quantization through mean-centered
  transformations.
\newblock In \emph{ES-FoMo III Workshop at the International Conference on
  Machine Learning}, 2025.
\newblock arXiv:2506.09932.

\bibitem{frantar2023}
E.~Frantar, S.~Ashkboos, T.~Hoefler, and D.~Alistarh.
\newblock {GPTQ}: Accurate post-training quantization for generative
  pre-trained transformers.
\newblock In \emph{International Conference on Learning Representations}, 2023.
\newblock arXiv:2210.17323.

\bibitem{he2025}
H.~He and Thinking Machines Lab.
\newblock Defeating nondeterminism in {LLM} inference.
\newblock \emph{Thinking Machines Lab: Connectionism}, 2025.
\newblock
  \url{https://thinkingmachines.ai/blog/defeating-nondeterminism-in-llm-inference/}.

\bibitem{maharana2024}
A.~Maharana, D.-H. Lee, S.~Tulyakov, M.~Bansal, F.~Barbieri, and Y.~Fang.
\newblock Evaluating very long-term conversational memory of {LLM} agents.
\newblock In \emph{Proceedings of ACL}, 2024.
\newblock arXiv:2402.17753.

\bibitem{reimers2019}
N.~Reimers and I.~Gurevych.
\newblock Sentence-{BERT}: Sentence embeddings using siamese {BERT}-networks.
\newblock In \emph{Proceedings of EMNLP-IJCNLP}, 2019.
\newblock arXiv:1908.10084.

\bibitem{thakur2021}
N.~Thakur, N.~Reimers, A.~R{\"u}ckl{\'e}, A.~Srivastava, and I.~Gurevych.
\newblock {BEIR}: A heterogeneous benchmark for zero-shot evaluation of
  information retrieval models.
\newblock In \emph{NeurIPS Datasets and Benchmarks Track}, 2021.
\newblock arXiv:2104.08663.

\bibitem{wang2020minilm}
W.~Wang, F.~Wei, L.~Dong, H.~Bao, N.~Yang, and M.~Zhou.
\newblock {MiniLM}: Deep self-attention distillation for task-agnostic
  compression of pre-trained transformers.
\newblock In \emph{Advances in Neural Information Processing Systems}, 2020.
\newblock arXiv:2002.10957.

\bibitem{wang2022e5}
L.~Wang, N.~Yang, X.~Huang, B.~Jiao, L.~Yang, D.~Jiang, R.~Majumder, and F.~Wei.
\newblock Text embeddings by weakly-supervised contrastive pre-training.
\newblock arXiv:2212.03533, 2022.

\bibitem{warner2025}
B.~Warner, A.~Chaffin, B.~Clavi{\'e}, O.~Weller, O.~Hallstr{\"o}m,
  S.~Taghadouini, A.~Gallagher, R.~Biswas, F.~Ladhak, T.~Aarsen, N.~Cooper,
  G.~Adams, J.~Howard, and I.~Poli.
\newblock Smarter, better, faster, longer: A modern bidirectional encoder for
  fast, memory efficient, and long context finetuning and inference.
\newblock In \emph{Proceedings of ACL}, 2025.
\newblock arXiv:2412.13663.

\bibitem{xiao2023}
G.~Xiao, J.~Lin, M.~Seznec, H.~Wu, J.~Demouth, and S.~Han.
\newblock {SmoothQuant}: Accurate and efficient post-training quantization for
  large language models.
\newblock In \emph{International Conference on Machine Learning}, 2023.
\newblock arXiv:2211.10438.

\bibitem{yvinec2023}
E.~Yvinec, A.~Dapogny, M.~Cord, and K.~Bailly.
\newblock {REx}: Data-free residual quantization error expansion.
\newblock In \emph{Advances in Neural Information Processing Systems}, 2023.
\newblock arXiv:2203.14645.

\bibitem{zhang2025sage}
J.~Zhang, J.~Wei, P.~Zhang, J.~Zhu, and J.~Chen.
\newblock {SageAttention}: Accurate 8-bit attention for plug-and-play inference
  acceleration.
\newblock In \emph{International Conference on Learning Representations}, 2025.
\newblock arXiv:2410.02367.

\bibitem{zhang2025sage2}
J.~Zhang, H.~Huang, P.~Zhang, J.~Wei, J.~Zhu, and J.~Chen.
\newblock {SageAttention2}: Efficient attention with thorough outlier smoothing
  and per-thread {INT4} quantization.
\newblock In \emph{International Conference on Machine Learning}, 2025.
\newblock arXiv:2411.10958.

\end{thebibliography}

\end{document}
